\phantomsection\label{fig:vol01:spring-autumn:boju}
\begin{center}
\begin{tikzpicture}[x=.88cm,y=.88cm,font=\small]
  \fill[paper] (0,0) rectangle (10.5,5.7);
  \draw[blue!50,line width=.9pt] (.4,4.8) .. controls (2.5,4.6) and (4.3,4.9) .. (6.3,4.4);
  \node[text=blue!70!black,above] at (3.5,4.75) {淮河通道};
  \draw[blue!55,line width=1pt] (5.2,.4) .. controls (6.1,1.8) and (7.0,2.2) .. (9.8,1.5);
  \node[text=blue!70!black,below] at (7.3,1.75) {汉水与江汉平原};

  \fill[cyan!22] (.8,3.2) -- (3.4,3.0) -- (3.8,4.8) -- (1.2,5.1) -- cycle;
  \node[align=center] at (2.1,4.0) {吴\\下游水网};
  \fill[purple!20] (5.4,1.1) -- (9.8,1.0) -- (10.0,4.4) -- (6.3,4.6) -- cycle;
  \node[align=center] at (7.8,3.0) {楚\\江汉腹地};
  \fill[orange!20] (8.7,4.4) -- (10.1,4.4) -- (10.2,5.4) -- (8.9,5.5) -- cycle;
  \node at (9.5,4.95) {郢都方向};
  \fill[timeline] (5.2,3.5) circle (.14);
  \node[above,align=center] at (5.2,3.6) {柏举\\战场};

  \draw[-{Stealth[length=2mm]},ultra thick,color=cyan!60!black] (2.7,3.8) -- (5.05,3.5);
  \draw[-{Stealth[length=2mm]},ultra thick,color=cyan!60!black] (5.35,3.45) -- (8.5,4.75);
  \node[text=cyan!60!black,above] at (4.0,3.65) {吴军经北线深入};
  \node[text=cyan!60!black,left] at (7.1,4.15) {向郢推进};
\end{tikzpicture}

\smallskip
\small\textit{图  柏举战役空间示意：强调吴军利用北方通道进入楚地、柏举战场与郢都方向的关系；不表示精确战线、河道或疆界。}
\end{center}
